\clearpage
\section{Experimental setup}\label{sec:setup}

All the measurements of this report were made on the computer described in Table~\ref{tab:setup}.
\begin{table}[H]
\centering
\begin{tabular}{@{}p{4.6cm}p{11.4cm}@{}}
\toprule
Component & Configuration\\
\midrule
Computer & Acer Nitro ANV16-42-R3JW\\
Processor & AMD Ryzen 7 260 with Radeon 780M Graphics, x86\_64\\
Cores & 1 socket, 8 physical cores, 2 threads per core, 16 logical CPUs\\
Cache & L1 data 256 KiB and L1 instruction 256 KiB (8 instances each), L2 8 MiB (8 instances), L3 16 MiB\\
CPU frequency policy & scaling governor \texttt{powersave}\\
Memory & 14 GiB of RAM (15\,126\,480 kB), 4 GiB of swap\\
Memory modules & 2 modules of 8 GB (16 GB), DDR5, speed 5600 MT/s\\
Graphics & NVIDIA GeForce RTX 5060 (dedicated, 8 GB GDDR7) and the integrated AMD Radeon 780M; the solver uses neither\\
Storage & NVMe drive; the file system holding the project has 128 GB\\
Operating system & Ubuntu 26.04.1 LTS, Linux kernel 7.0.0-38-generic\\
Python & CPython 3.14.4\\
\bottomrule
\end{tabular}
\caption{Configuration of the computer used for the measurements.}
\label{tab:setup}
\end{table}

The measurements of \hyperref[sec:q2]{Question~2}, \hyperref[sec:q5]{Question~5} and \hyperref[sec:q6]{Question~6} ran several solver processes at the same time, eight, ten and eight respectively, on the 16 logical CPUs; those of \hyperref[sec:q3]{Question~3} ran one process at a time. The processes that run together share the caches and the memory bandwidth, so the times are indicative.

This computer is not the one of the final competition. The assignment specifies an Intel Core Ultra 7 155H, Ubuntu 22.04 and Python 3.12 in a Docker container limited to 4 CPUs and 8 GiB of memory, with a limit of 60 seconds per instance. The times of this report cannot be compared directly with those of the competition.
